\chapter{Wi-Fi: a pure-Ada driver over a binary radio}\index{Wi-Fi}\index{WiFi}\index{802.11}
\label{ch:wifi}

The networking chapter (\ref{ch:networking}) gave the application a wired NIC and
the standard \texttt{GNAT.Sockets} API; the TLS chapter (\ref{ch:tls}) put HTTPS
over it. This chapter gives a \emph{wireless} NIC --- the ESP32-S3's on-chip
radio --- and the reward is almost anticlimactic: the same socket stack and the
same TLS client run over it unchanged. The \texttt{esp32s3\_wifi\_tls} example is
the wired \texttt{esp32s3\_tls\_weather} with one line swapped underneath.

But Wi-Fi is unlike every other driver in this book in one honest way, and it is
worth stating plainly up front. Every peripheral so far --- SPI, I2C, the W5500,
the SD card, the crypto accelerators --- we drove ourselves, register by
register, in Ada. The 802.11 \emph{lower MAC} (the microsecond-scale timing of
acknowledgements, back-off and rate control) and the \emph{PHY} (RF calibration
for a specific silicon revision) are not ours to write: Espressif ships them only
as binaries, and there is no public source to reimplement from. So this is the
one driver in the book that links a blob. The interesting engineering is what
surrounds that blob --- all of which \emph{is} ours, and all of which is Ada.

\cnote{On ESP-IDF, Wi-Fi is \texttt{esp\_wifi} standing on these same blobs,
scheduled by FreeRTOS and feeding lwIP's BSD sockets. We keep the blobs (there is
no alternative) but throw the rest away: no FreeRTOS, no lwIP. The blob's RTOS
calls land on our Jorvik tasks, and its frames land on the pure-Ada stack from
Chapter~\ref{ch:networking}.}

\begin{figure}[ht]
\centering
\begin{tikzpicture}[node distance=2mm]
\node[layer, fill=blue!9] (app)  {Your application: portable \texttt{GNAT.Sockets} + \texttt{DNS\_Client} + \texttt{TLS\_Client}};
\node[layer, below=of app]  (fac)  {\texttt{GNAT.Sockets} facade / \texttt{Net\_Devices.Device'Class} (Chapters \ref{ch:networking}, \ref{ch:tls})};
\node[layer, below=of fac]  (nd)   {\texttt{ESP32S3.WiFi.Net\_Device}: the radio as a NIC backend};
\node[layer, below=of nd]   (ip)   {\texttt{ESP32S3.WiFi.IP}: pure-Ada Ethernet/ARP/IPv4/UDP/TCP};
\node[layer, below=of ip]   (drv)  {\texttt{ESP32S3.WiFi} + \texttt{.Supplicant} + \texttt{.Port}: bring-up, WPA2, tx/rx};
\node[layer, below=of drv, fill=red!8]  (blob) {\textbf{Espressif blobs}: \texttt{libnet80211}, \texttt{libpp}, \texttt{libcore}, \texttt{libphy}};
\node[layer, below=of blob] (os)   {\texttt{ESP32S3.WiFi.RTOS}: the OS adapter --- the blob's FreeRTOS, re-implemented on Jorvik};
\end{tikzpicture}
\caption{The Wi-Fi stack. Everything except the one red layer is Ada. Above
\texttt{Net\_Devices.Device'Class} nothing changes from the wired chapters; the
blob is boxed in between a thin Ada driver above it and an Ada OS adapter beneath
it.}
\label{fig:wifi-layers}
\end{figure}

\section{The one place we link a blob}\index{Wi-Fi!blobs}\index{blob!fetching}
\label{ch:wifi:blobs}

Four archives make up the radio firmware, pinned to ESP-IDF v5.4.4:
\texttt{libnet80211.a} (the MLME: scanning, association, the state machine),
\texttt{libpp.a} (the lower-MAC ``packet processor''), \texttt{libcore.a} (the
\texttt{esp\_wifi} entry points) and \texttt{libphy.a} (RF calibration). Together
they are about 2.3\,MB.

We do not commit them. A repository of pure-Ada source should not carry megabytes
of someone else's binaries, and it does not need to: both upstream libraries
(\texttt{espressif/esp32-wifi-lib} and \texttt{espressif/esp-phy-lib}) are
\textbf{Apache-2.0} licensed and live at stable Git commits. So the SDK ships a
manifest and a fetch script instead:

\begin{lstlisting}[language=bash]
$ tools/fetch-wifi-blobs.sh
fetch  libnet80211.a  <-  espressif/esp32-wifi-lib@7bd1599
  verified + installed
...
\end{lstlisting}

The script reads \texttt{libs/esp32s3\_wifi/blobs/MANIFEST.lock}, downloads each
archive from its pinned commit, and verifies it against a recorded SHA-256 before
installing it. The \texttt{.a} files are \texttt{.gitignore}d; only the manifest,
the two Apache-2.0 licence texts and a provenance note are tracked. A build with
no local ESP-IDF fetches once and caches; a developer who \emph{does} have
ESP-IDF sets \texttt{IDF\_PATH} and links their own copies. Either way the result
is bit-identical --- that is what the checksum guarantees.

\cnote{This is the standard ``vendored binary'' pattern: pin an exact upstream
commit, record a checksum, fetch on demand. It keeps history free of binaries
\emph{and} makes the build reproducible and tamper-evident --- an unpinned
``download the latest'' would be neither.}

\section{Teaching a FreeRTOS blob to live on Jorvik}\index{Wi-Fi!OS adapter}\index{FreeRTOS}
\label{ch:wifi:osadapter}

The blob was compiled expecting an operating system. It calls \texttt{task\_create}
to spawn its worker; it pushes received frames through \texttt{queue\_send} and
blocks the worker on \texttt{queue\_recv}; it takes recursive mutexes around MAC
state; it arms one-shot timers for retransmission; it calls \texttt{malloc} for
frame buffers. On a stock system those are FreeRTOS. Here they are a few hundred
lines of Ada --- the \emph{OS adapter}, \texttt{ESP32S3.WiFi.RTOS} --- and every
one of them is a Jorvik construct in disguise.

The blob's ``tasks'' are the clearest example. Ravenscar forbids creating tasks
dynamically, so instead of honouring \texttt{task\_create} literally we hand the
blob one of a fixed pool of \emph{worker} tasks, each parked on a mailbox waiting
for a procedure to run:

\begin{lstlisting}
task type Worker with Priority => Worker_Prio,
                      CPU      => Wifi_CPU,      --  pin to the ISR's core
                      Storage_Size => Worker_Stk;

task body Worker is
   Idx  : Natural;
   F, P : System.Address;
begin
   Allocator.Claim_Worker (Idx);          --  take a pool slot
   loop
      Mailboxes (Idx).Wait (F, P);        --  block until task_create hands us one
      if F /= System.Null_Address then
         To_Entry (F) (P);                --  run the blob's "task body" -- forever
      end if;
   end loop;
end Worker;
\end{lstlisting}

\texttt{task\_create} finds a free worker, posts the blob's entry point to its
mailbox, and returns the mailbox address as the ``task handle'' --- which matters,
because the blob later compares handles to decide whether it is running
\emph{on} the Wi-Fi task, and the OS adapter must give back exactly what it
handed out. Queues become protected objects with a bounded ring and a
\texttt{when Count > 0} barrier; a blocking \texttt{queue\_recv} is a Jorvik
entry call, an interrupt-context \texttt{queue\_send\_from\_isr} is a protected
procedure. Mutexes are protected objects with an ownership count. Timers are a
single high-priority \texttt{Timer\_Task} servicing a small table. \texttt{malloc}
maps onto the DMA-capable internal-RAM heap. The blob thinks it is on FreeRTOS;
FreeRTOS is not linked.

\cnote{A recursive mutex over 802.11 state, taken and released thousands of times
a second, must be genuinely cheap. A Jorvik protected object is: enter, touch a
field, exit --- no kernel call, priority-ceiling locking on a single core. The
adapter is not a compatibility shim bolted on the side; it \emph{is} the RTOS the
blob runs on.}

\section{Bring-up: from reset to a scan}\index{Wi-Fi!bring-up}\index{PHY}\index{WPA2}
\label{ch:wifi:bringup}

With the OS adapter in place, bring-up follows the sequence the blob expects:
initialise the \texttt{esp\_wifi} core, run PHY calibration (the reason
\texttt{libphy} exists --- the radio must be tuned to \emph{this} chip and
antenna), start the MAC, and route the WMAC interrupt. That last step is pure
Ada: the lower MAC raises a hardware interrupt when a frame arrives, and GNARL
delivers it to a protected handler like any other device interrupt.

\begin{lstlisting}
protected WMAC_ISR
  with Interrupt_Priority => Ada.Interrupts.Names.Device_L2_Priority
is
   procedure Handler
     with Attach_Handler => Ada.Interrupts.Names.Device_L2_1;
end WMAC_ISR;

protected body WMAC_ISR is
   procedure Handler is
   begin
      if Blob_Isr /= System.Null_Address then
         To_Isr (Blob_Isr) (Blob_Arg);   --  hand the frame to the blob's ISR
      end if;
   end Handler;
end WMAC_ISR;
\end{lstlisting}

The blob's ISR is closure-free and library-level, exactly as
Chapter~\ref{ch:interrupts} requires --- a captured environment here would fault
on the trampoline the ESP32-S3 cannot execute.

PHY calibration is also where the radio's \emph{power management} lives. Once the
core is up, the lower MAC does not keep the RF lit continuously: it drives the
\texttt{phy\_enable} / \texttt{phy\_disable} hooks roughly once a second, closing
the analog front-end whenever it has nothing to send or receive. The adapter
honours that cycle exactly as ESP-IDF does on this chip --- \texttt{phy\_disable}
saves the PHY digital registers (\texttt{phy\_dig\_reg\_backup}), calls
\texttt{phy\_close\_rf} and powers down the temperature sensor; \texttt{phy\_enable}
re-wakes the PHY with \texttt{phy\_wakeup\_init} and restores those registers. The
save/restore is not optional on the ESP32-S3: without it the RF degrades over
repeated close/wake cycles. The upshot is that the radio powers down between
activity on its own, which is the groundwork for the low-power modes of
Chapter~\ref{ch:lowpower}; the deep-sleep path there does not even need it, since
the RTC power-down takes the whole RF domain regardless.

Scanning drops out of this for free: \texttt{ESP32S3.WiFi.Scan} asks the MLME to
sweep the channels and reads back the results. Connecting adds the
\emph{supplicant} --- the WPA2 four-way handshake in \texttt{ESP32S3.WiFi.Supplicant}
--- which derives the pairwise key and installs it into the hardware AES-CCMP
engine so that frames are decrypted in silicon. (Getting that key to the right
slot of the crypto engine was, for the record, the single most stubborn bug in
the whole bring-up; JTAG earned its keep.)

\section{Wi-Fi as just another NIC}\index{Wi-Fi!Net\_Device}\index{Net\_Devices}
\label{ch:wifi:netdevice}

Here is the payoff for all the interface discipline of
Chapter~\ref{ch:networking}. The blob hands us decrypted Ethernet frames and
accepts frames to transmit; that is precisely a NIC. So
\texttt{ESP32S3.WiFi.Net\_Device} implements the same
\texttt{Net\_Devices.Device'Class} the W5500 does, and feeds a small pure-Ada IP
stack (\texttt{ESP32S3.WiFi.IP}: Ethernet framing, ARP, IPv4, UDP and a
stop-and-wait TCP). Above the \texttt{Device'Class} line, \emph{nothing knows it
is talking to a radio}:

\begin{lstlisting}
--  esp32s3_wifi_tls, after association + DHCP:
Session : TLS_Client.Session;
...
Connect (Wifi_Credentials.SSID, Wifi_Credentials.Pass, Result => St);
...  --  DHCP, then DNS, then:
TLS_Client.Handshake (Session, Host => "api.open-meteo.com", Port => 443);
\end{lstlisting}

That \texttt{TLS\_Client}, that \texttt{DNS\_Client}, that TCP --- all of it is
the code from the previous two chapters, unmodified. The full pipeline runs on
hardware against a live access point and the real \texttt{api.open-meteo.com}:
associate, DHCP-lease an address, resolve the name, fetch the time over NTP (TLS
needs a wall clock to judge certificate validity), complete a TLS 1.3 handshake,
RSA-PSS-verify the server's \texttt{CertificateVerify}, validate the chain up to
the real ISRG Root X1, and GET the forecast. Pure-Ada HTTPS, over a pure-Ada
stack, over a radio we do not have the source to.

\section{A war story: the system that slept forever}\index{Wi-Fi!deadlock}\index{waiti}\index{CCOUNT}\index{SYSTIMER}
\label{ch:wifi:deadlock}

Wi-Fi is the reason a subtle runtime bug finally surfaced, and the hunt is worth
telling because it is exactly the kind of thing bare-metal throws at you.

The TLS-over-Wi-Fi example built cleanly and then hung during boot --- before its
first line of output. No crash, no exception: a \emph{silent} hang, the worst
kind. Under OpenOCD and GDB, \texttt{info tasks} told the story at a glance: the
environment task and the Wi-Fi timer task were both in \texttt{Delay Sleep}, the
worker tasks were idle on their mailboxes, and \emph{both} CPU cores were sitting
in \texttt{waiti}. Nothing was running, nothing was ready, and nothing was going
to change that. Reading the cycle counter twice, a second apart, gave the
diagnosis: \texttt{CCOUNT} was frozen.

The runtime clocked its alarm off the Xtensa \texttt{CCOUNT}/\texttt{CCOMPARE2}
timer --- and \texttt{CCOUNT stops advancing while the core sits in
\texttt{waiti}}. As long as \emph{something} ran, that was invisible. But the
Wi-Fi example, during elaboration, reached a moment where every task was blocked
on a timed delay before any device interrupt existed to wake a core. Both cores
idled, \texttt{CCOUNT} froze, the alarm deadline was never reached, and the only
thing that could have re-armed the alarm was the alarm itself. A fully-idle
system could never wake. (The same latent flaw had been mis-diagnosed for months
as delays occasionally waking ``about fifteen seconds late'' --- that fifteen
seconds is exactly how long \texttt{CCOUNT} takes to wrap all the way around.)

The fix moves the alarm off \texttt{CCOUNT} and onto the \emph{SYSTIMER}, a
free-running 16\,MHz counter that keeps ticking through \texttt{waiti}. Each core
arms one of its comparators; the comparator's interrupt is routed to a level-5
CPU interrupt, so the existing vector and dispatch (Chapter~\ref{ch:interrupts})
are untouched --- only the timer underneath changed:

\begin{lstlisting}
procedure Set_Alarm (Ticks : Timer_Interval) is
   Core     : constant Integer := Integer (Current_CPU) - 1;
   Now      : constant Unsigned_64 := Native_Systimer_Count;   --  free-running
   Deadline : constant Unsigned_64 := Now + Unsigned_64'Max (1, U64 (Ticks) / 15);
begin
   --  Arm this core's SYSTIMER comparator (one-shot); it fires on
   --  count >= target, so even a past deadline fires at once -- no exact-match
   --  miss, unlike CCOMPARE2.
   Arm_Comparator (Core, Deadline);
end Set_Alarm;
\end{lstlisting}

It is a runtime-wide change --- it is now the timer behind \emph{every}
\texttt{delay until} in the system --- so it was validated well beyond Wi-Fi: an
SMP example, and a dedicated \texttt{esp32s3\_delay\_test} that measures the
error of a delay taken from a fully-idle system and finds it exact. The lesson is
an old one, sharpened: a clock that stops when the CPU rests is not a clock you
can set an alarm by.

\section{The examples}\index{Wi-Fi!examples}
\label{ch:wifi:examples}

\begin{itemize}[nosep]
\item \texttt{esp32s3\_wifi\_scan} --- bring up the radio and list nearby access
points with channel, RSSI and BSSID.
\item \texttt{esp32s3\_wifi\_sniff} --- promiscuous-mode 802.11 capture, the MAC
handed raw frames off the air.
\item \texttt{esp32s3\_wifi\_dns} --- associate and resolve a name through
\texttt{DNS\_Client}.
\item \texttt{esp32s3\_wifi\_http} --- a plain HTTP GET over the software TCP
stack.
\item \texttt{esp32s3\_wifi\_tls} --- the capstone: full TLS 1.3 HTTPS end to end,
described above.
\item \texttt{esp32s3\_wifi\_deepsleep} --- associate, then deep-sleep: entering
deep sleep is itself the radio power-down (the RTC cuts the RF domain), and the
stack re-initialises on each wake (Chapter~\ref{ch:lowpower}).
\end{itemize}

Each needs the blobs (\texttt{tools/fetch-wifi-blobs.sh}, or \texttt{IDF\_PATH})
and your network in the git-ignored \texttt{src/wifi\_credentials.ads} --- copy
the \texttt{.template} and fill in the SSID and passphrase. The radio is not
ours; everything you can read is.
